\section*{Chapitre \ref{ch:b3:sensory-systems} --- Les systèmes sensoriels}

\begin{solution}{exo:b3:sensory-systems:1}
Par le fil, non par le message~: chaque nerf sensoriel est une \omterm{def:b3:sensory-systems:principles}{ligne
étiquetée} qui aboutit dans sa propre région du cerveau, et tout ce qui
arrive par le nerf optique est lu comme de la lumière --- et c'est
pourquoi une pression sur le globe oculaire produit un éclair et un coup
sur l'oreille un tintement.
\end{solution}

\begin{solution}{exo:b3:sensory-systems:2}
Dans le noir, le GMP cyclique tient ouverts des canaux cationiques et un
«~courant d'obscurité~» entrant constant garde la cellule dépolarisée et
libérant du glutamate. La lumière active la \omterm{def:b3:sensory-systems:retina}{rhodopsine}, la transducine
et la phosphodiestérase, qui détruit le GMPc~; les canaux se ferment, le
courant entrant s'arrête, et la cellule s'hyperpolarise et libère moins
--- la lumière est signalée par une diminution.
\end{solution}

\begin{solution}{exo:b3:sensory-systems:3}
$I = I_{0}10^{L/10}$~: \qty{20}{dB}, \qty{1e-10}{W/m^2}~; \qty{60}{dB},
\qty{1e-6}{W/m^2}~; \qty{110}{dB}, \qty{0.1}{W/m^2}. Du plus fort au
plus faible~: $10^{9}$.
\end{solution}

\begin{solution}{exo:b3:sensory-systems:4}
Le sucré, l'umami et l'amer par des récepteurs couplés aux protéines~G
(un du sucré, un de l'umami, environ vingt-cinq de l'amer)~; le salé par
le canal sodique ENaC~; l'acide par un canal à protons. L'essentiel de
la «~saveur~» est l'odeur des molécules volatiles qui gagnent le nez
depuis la bouche~; le nez bouché, il ne reste que les cinq \omterm{def:b3:sensory-systems:chemical}{goûts} et les
aliments semblent fades.
\end{solution}

\begin{solution}{exo:b3:sensory-systems:5}
Fechner~: $S = (1/0.02)\ln(10\,000/100) = 50\ln 100 \approx 230$ pas. En
comptant~: $1.02^{n} = 100$ donne $n = \ln 100/\ln 1.02 = 233$.
\end{solution}

\begin{solution}{exo:b3:sensory-systems:6}
$P = 1 - e^{-a}\sum_{k=0}^{5}a^{k}/k!$~: $a = 2$~: $0.017$~; $a = 6$~:
$0.55$~; $a = 12$~: $0.98$. La moitié des éclairs sont vus à
$a \approx 5.7$~; avec \qty{10}{\%} d'absorption l'éclair doit délivrer
environ $57$ photons à la cornée.
\end{solution}

\begin{solution}{exo:b3:sensory-systems:7}
$55/3\times 1.3 \approx 24$~; en \omterm{def:b3:sensory-systems:cochlea}{décibels} de pression, $20\log_{10}24
\approx \qty{28}{dB}$. Des osselets soudés ne peuvent transmettre la
vibration, si bien que le son n'atteint la \omterm{def:b3:sensory-systems:cochlea}{cochlée} que par conduction à
travers le crâne, sans le gain de l'oreille moyenne~: une perte de
transmission de quelque \qtyrange{30}{50}{dB}.
\end{solution}

\begin{solution}{exo:b3:sensory-systems:8}
Le canal de la \omterm{def:b3:sensory-systems:cochlea}{cellule ciliée} est ouvert directement par le \omterm{def:b3:sensory-systems:cochlea}{lien apical}
qui tire dessus --- pas de messager, pas d'enzyme --- de sorte qu'il
répond en microsecondes~; le photorécepteur amplifie un photon par deux
étages enzymatiques, qui prennent des dizaines de millisecondes.
L'audition gagne la précision temporelle qu'il faut pour localiser les
sons par les délais interauraux et pour suivre la phase~; la vision
gagne la sensibilité qui permet de compter des photons uniques, et le
paie en rapidité.
\end{solution}

\begin{solution}{exo:b3:sensory-systems:9}
$H(0.05) = -0.05\ln 0.05 - 0.95\ln 0.95 = 0.150 + 0.049 = 0.199$.
Humain~: $\ln\binom{400}{20} \approx 400\times 0.199 = 80$, environ
$10^{34}$ motifs. Souris~: $\ln\binom{1100}{55} \approx 1100\times 0.199
= 219$, environ $10^{95}$.
\end{solution}

\begin{solution}{exo:b3:sensory-systems:10}
Moyenne $500\times 0.1/40 = 1.25$ événement spontané par fenêtre. $P(k
\le 5) = e^{-1.25}(1 + 1.25 + 0.78 + 0.33 + 0.10 + 0.03) = 0.998$, donc
$P(k \ge 6) \approx 2\times 10^{-3}$. Avec un seuil de un, un événement
spontané se produirait dans la plupart des fenêtres
($1 - e^{-1.25} = 0.71$) et le noir serait plein d'éclairs~; à six, les
fausses alarmes viennent une fois sur cinq cents fenêtres. Le seuil est
fixé là où le bruit des \omterm{def:b3:sensory-systems:retina}{bâtonnets}, et non leur sensibilité, le dicte.
\end{solution}

\begin{solution}{exo:b3:sensory-systems:11}
Plus grand délai $0.2/340 = \qty{590}{\micro s}$ (son venant d'un côté).
Près du plan médian $\Delta t \approx (d/c)\sin\theta$, de sorte que
$\qty{10}{\micro
s}$ correspond à $\sin\theta = 0.017$, environ \qty{1}{\degree}.
Au-dessus de \qty{4}{kHz} la période (\qty{250}{\micro s}) est plus
courte que ce à quoi les neurones peuvent se verrouiller, et un délai de
plusieurs centaines de microsecondes couvre plus d'un cycle, ce qui rend
la phase ambiguë~; le cerveau emploie alors l'écart de niveau entre les
oreilles et le code de place.
\end{solution}

\begin{solution}{exo:b3:sensory-systems:12}
Loin du bord, $r = I(1 - 2w)$~: $0.6$ du côté sombre, $1.2$ du côté
clair. Au dernier récepteur sombre ($I = 1$, voisins $1$ et $2$)~: $r
= 1 - 0.2\times 3 = 0.4$~; au premier clair~: $r = 2 - 0.6 = 1.4$. La
marche est exagérée en un creux et une bosse --- les bandes de Mach. La
\omterm{def:b3:sensory-systems:retina}{rétine} y gagne le contraste et les bords, et une dynamique moindre à
transmettre~; elle y perd la fidélité aux niveaux absolus, et produit
des illusions de clarté.
\end{solution}

\begin{solution}{pb:b3:sensory-systems:1}
\textbf{1.} $0.01/3.6\times 10^{-19} = 2.8\times 10^{16}$ photons par
seconde.
\textbf{2.} Aire de la pupille $\pi(3.5\times 10^{-3})^{2} = 3.8\times 10^{-5}$
m$^{2}$. À \qty{1}{km}~: fraction $3.8\times 10^{-5}/1.26\times 10^{7} =
3\times 10^{-12}$, $8.5\times 10^{4}$ photons par seconde. À
\qty{10}{km}~: $850$ par seconde.
\textbf{3.} En \qty{0.1}{s} avec \qty{10}{\%} d'absorption~: $850$ à
\qty{1}{km}, $8.5$ à \qty{10}{km} --- tous deux au-dessus de six~; la
bougie est vue à dix kilomètres (la plupart du temps).
\textbf{4.} Six absorbés par fenêtre demandent $600$ photons par seconde
dans la pupille~: $4\pi d^{2} = 3.8\times 10^{-5}\times 2.8\times 10^{16}/600
= 1.8\times 10^{9}$ m$^{2}$, $d \approx \qty{12}{km}$. Là, $P(k \ge
6 \mid a = 6) = 0.55$~: la bougie est vue dans environ la moitié des
fenêtres.
\textbf{5.} $e^{-6} = 0.0025$. Au seuil le compte fluctue d'une fenêtre
à l'autre --- $6 \pm 2.4$ --- de sorte que la source semble apparaître et
disparaître~: le scintillement d'une étoile faible (auquel l'atmosphère
ajoute le sien).
\textbf{6.} La fluctuation du fond, $\sqrt{10^{4}} = 100$ photons par
fenêtre, noie les $8.5$ de la bougie~: la détection demande que le
signal dépasse le bruit du fond, non le seuil de l'œil dans le noir.
\textbf{7.} $\qty{1}{pA}/200 = \qty{5}{fA}$ par canal~; $5\times
10^{-15}/1.6\times 10^{-19} = 3\times 10^{4}$ ions par seconde.
\textbf{8.} Un trentième~; une trentaine de photons simultanés ferment
tous les canaux et saturent le \omterm{def:b3:sensory-systems:retina}{bâtonnet}.
\textbf{9.} $10^{-12}\times 0.2 = 2\times 10^{-13}$ C, $1.3\times 10^{6}$
cations~: un photon commande plus d'un million de charges.
\textbf{10.} Le \omterm{def:b3:sensory-systems:retina}{bâtonnet} sature en quelques millisecondes, tous canaux
fermés, et sa \omterm{def:b3:sensory-systems:retina}{rhodopsine} est blanchie plus vite qu'elle n'est
régénérée~; les \omterm{def:b3:sensory-systems:retina}{cônes}, à gain plus faible et extinction plus rapide,
continuent de répondre et déplacent leur plage sur six décades de
lumière.
\textbf{11.} Un gain élevé demande une longue intégration (chaque étage
prend du temps et la réponse est prolongée pour s'accumuler), de sorte
que le \omterm{def:b3:sensory-systems:retina}{bâtonnet} est lent et sensible~; l'amplification moindre du \omterm{def:b3:sensory-systems:retina}{cône}
et son extinction plus prompte donnent une réponse en \qty{20}{ms},
assez vite pour le mouvement, au prix d'une centaine de photons
nécessaires.
\textbf{12.} $10^{8}/40 = 2.5\times 10^{6}$ isomérisations spontanées par
seconde sur toute la \omterm{def:b3:sensory-systems:retina}{rétine} --- mais seulement $1.25$ par fenêtre dans
une plage de $500$ \omterm{def:b3:sensory-systems:retina}{bâtonnets}, bien au-dessous du seuil de six, de sorte
que le bruit est rejeté plage par plage~; le seuil et la mise en commun
existent précisément pour cela.
\textbf{13.} \qty{0.1}{W/m^2}~; sur \qty{5.5e-5}{m^2}, \qty{5.5e-6}{W}.
\textbf{14.} $5.5\times 10^{-6}\times 7200 = \qty{0.04}{J}$ --- environ
quatre cents gouttes de pluie en deux heures.
\textbf{15.} $10^{-12}\times 5.5\times 10^{-5} = 5.5\times 10^{-17}$ W~;
en \qty{10}{ms}, $5.5\times 10^{-19}$ J --- à peu près l'énergie d'un ou
deux photons visibles.
\textbf{16.} $0.3/5000 = 6\times 10^{-5}$ rad, $0.003^{\circ}$~; la
déflexion vaut environ trois atomes d'hydrogène.
\textbf{17.} $10^{(110 - 85)/10} = 316$~: deux heures à \qty{110}{dB}
délivrent l'énergie de $630$ heures à \qty{85}{dB} --- près d'un mois de
jours ouvrés.
\textbf{18.} Environ $120$ pas juste perceptibles, un par \omterm{def:b3:sensory-systems:cochlea}{décibel}.
\textbf{19.} $2^{400} = 10^{400\log_{10}2} = 10^{120}$.
\textbf{20.} $\ln\binom{400}{20} \approx 400\,H(0.05) \approx 80$, moins
une correction de Stirling d'environ $2$~: quelque $10^{34}$ motifs.
\textbf{21.} Compter les récepteurs présents dans un motif et pas dans
l'autre~: $5 +
5 = 10$ sur $40$ (une distance de Hamming)~; des motifs qui se
recouvrent aux trois quarts activent presque les mêmes \omterm{def:b3:sensory-systems:chemical}{glomérules} et
sont lus comme presque la même odeur.
\textbf{22.} $10^{12}/10^{34} = 10^{-22}$ des motifs. Le code n'est pas
la limite~: les récepteurs sont bruités, beaucoup sont accordés de la
même façon et leurs activités sont donc corrélées, les molécules réelles
ne produisent qu'un petit sous-ensemble de motifs, et la discrimination
et la mémoire des motifs par le cerveau sont finies.
\textbf{23.} $\ln\binom{1100}{20} \approx 1100\,H(0.018) = 1100\times
0.091 \approx 100$, environ $10^{43}$ --- quelque $10^{9}$ fois le
nombre humain de motifs à vingt récepteurs (et bien davantage si la
souris emploie plus de récepteurs par odeur).
\textbf{24.} L'odorat général ne change guère --- le code est redondant
et d'autres récepteurs portent chaque odorant --- mais un odorant qui
dépend de ce récepteur à faible concentration devient indétectable ou
change de caractère~: une anosmie spécifique, comme chez les nombreuses
personnes qui ne sentent pas l'androsténone.
\textbf{25.} La bougie est vue jusqu'à environ \qty{12}{km}~; le son au
seuil délivre $5.5\times 10^{-19}$ J au tympan en \qty{10}{ms}, un
photon ou deux~; un nez humain peut former quelque $10^{34}$ motifs à
vingt récepteurs.
\end{solution}
